% Revised from the Aug18 manuscript using the October 1, 2026 replacements.
% Integration date: October 8, 2026.

\documentclass[10pt,reqno]{amsart}
% Unused bbm package omitted in standalone rendering.
\usepackage{mathrsfs}
\usepackage{amsfonts} %%% i.e. use 12pt type
\usepackage[dvipsnames,usenames]{color}
\textwidth=13.5cm %%% in the preamble; this will require
\baselineskip=17pt %%% after \begin{document}
\usepackage{graphicx,latexsym,bm,amsmath,amssymb,verbatim,multicol,lscape}
\vfuzz2pt % Don't report over-full v-boxes if over-edge is small
\hfuzz2pt % Don't report over-full h-boxes if over-edge is small
\newtheorem{thm}{Theorem} [section]
\newtheorem{cor}[thm]{Corollary}
\newtheorem{lem}[thm]{Lemma}
\newtheorem{exm}[thm]{Example}
\newtheorem{prop}[thm]{Proposition}
\theoremstyle{definition}
\newtheorem{defn}[thm]{Definition}
\theoremstyle{remark}
\newtheorem{rem}[thm]{Remark}
\newtheorem{con}[thm]{Conjecture}
\newtheorem{clm}[thm]{Claim}
\numberwithin{equation}{section}

\newcommand{\norm}[1]{\left\Vert#1\right\Vert}
\newcommand{\abs}[1]{\left\vert#1\right\vert}
\newcommand{\set}[1]{\left\{#1\right\}}
\newcommand{\Real}{\mathbb R}
\newcommand{\eps}{\varepsilon}
\newcommand{\To}{\longrightarrow}
\newcommand{\BX}{\mathbf{B}(X)}
\newcommand{\A}{\mathcal{A}}
\newcommand{\SSS}{\stackrel}
\newcommand{\V}{\mathrm{v}}
\newcommand{\N}{\mathbb N}
\newcommand{\Z}{\mathbb Z}
\newcommand{\R}{\mathbb R}
\newcommand{\Q}{\mathbb Q}
\newcommand{\F}{\mathcal F}
\newcommand{\C}{\mathbb C}
\newcommand{\MG}{\mathsf{MG}}
\newcommand{\rad}{\mathsf{rad}}
\newcommand{\sN}{\mathsf{N}}
\renewcommand{\subjclassname}
{\textup{2020} Mathematics Subject Classification}

\DeclareMathOperator{\ev}{ev}
\DeclareMathOperator{\ind}{ind}
\DeclareMathOperator{\res}{Res}
\DeclareMathOperator{\ord}{ord}
\DeclareMathOperator{\lcm}{lcm}
\DeclareMathOperator{\sgn}{sgn}
\DeclareMathOperator{\coker}{Coker}
\DeclareMathOperator{\spec}{spec}
\DeclareMathOperator{\supp}{supp}
\DeclareMathOperator{\Pic}{Pic}
\DeclareMathOperator{\cl}{cl}
\DeclareMathOperator{\Ker}{Ker}
\DeclareMathOperator{\Hom}{Hom}
\DeclareMathOperator{\st}{st}

\begin{document}
\title[Zero-sum sequences and additive bases]{Additive bases and the zero-sum invariant $\mathsf{N}(G)$ for groups $C_p\oplus C_{np}$}

\author[W.K. Yang]{Wenkai Yang}
\address{Center for Combinatorics, Nankai University, Tianjin 300071, P.R. China}
\email{3016210052@tju.edu.cn; 9820250116@nankai.edu.cn}

\thanks{This work was supported by
National Natural Science Foundation of China (Grant No. 12071344).}
\keywords{finite abelian group, additive basis, regular sequence, zero-sum subsequence}
\subjclass[2020]{Primary 11B75; Secondary 11A07, 20K01}

\begin{abstract}
Let $G$ be a finite abelian group, and let $c_0(G)$ be the least integer
such that every regular sequence over $G$ of at least this length is an
additive basis. Let $\mathsf N(G)$ be the least integer $l$ such that every
weak-regular sequence $S$ over $G$ with $|S|\ge l$ has a nonempty zero-sum
subsequence $T$ satisfying $\mathsf v_g(T)=\mathsf v_g(S)$ for some $g\mid S$.
We obtain a common upper bound for $c_0(C_q\oplus C_{nq})$ for groups
of odd order, where $q$ is prime, and prove $\mathsf N(G)\ge r(n_r-1)+1$ when
$G=C_{n_1}\oplus\cdots\oplus C_{n_r}$ with $1<n_1\mid\cdots\mid n_r$.
If $n\ge1$ and $p\ge3$ is the smallest prime divisor of $np$, we obtain
$\mathsf N(C_p\oplus C_{np^2})=2np^2-1$ and give upper bounds for
$\mathsf N(C_p\oplus C_{np})$ when $p\nmid n$.
We also determine $\mathsf N(C_n)$ for two infinite families of cyclic
groups arising in earlier work.
\end{abstract}

\maketitle

\section{Introduction and main results}

Let $G$ be a finite abelian group, with additive notation, and let $S$ be
a sequence over $G$. Denote by $\sum(S)$ the set of sums of nonempty
subsequences of $S$. The sequence $S$ is an additive basis of $G$ if
$\sum(S)=G$. For a subgroup $H\le G$, let $S_H$ be the subsequence
consisting of all terms of $S$ in $H$. We call $S$ regular over $G$ if
$|S_H|\le |H|-1$ for every proper subgroup $H<G$.
Let $c_0(G)$ be the least integer $l$ such that every regular sequence
over $G$ of length at least $l$ is an additive basis.
The problem of determining $c_0(G)$ was proposed by Olson;
see \cite{[GHQQZ2015],[GQZ2020]}.

Let $p$ be the smallest prime divisor of $|G|$, and let
\[
m(G)=
\begin{cases}
|G|,&G\text{ is cyclic},\\
2p-1,&G=C_p^2,\\
np+2p-3,&G=C_p\oplus C_{np},\ n\ge2,\\
|G|/p+p-2,&\text{otherwise}.
\end{cases}
\]
The inequality $c_0(G)\ge m(G)$ is known for every finite abelian group
\cite[Lemma~2.3]{[GQZ2020]}. Equality has been established for many
groups of rank at least three; see \cite{[GHQQZ2015],[GQZ2020]}.
For a group $G=C_{n_1}\oplus C_{n_2}$ with $1<n_1\mid n_2$,
the known cases include the following:
\begin{enumerate}
\item $|G|$ is even, or $G$ is a $p$-group
      \cite{[GHQQZ2015],[QH2017]};
\item $G=C_3\oplus C_{3q}$, where $q\ge5$ is prime
      \cite{[QL2022]};
\item $p\ge3$, $n_1\ge2p$ and $|G|\ge72p^6$
      \cite{[GLQW2022]};
\item $p\ge11$ and $n_1\ge p^2$
      \cite{[QLWZ2025]}.
\end{enumerate}
Qu, Li, Wang and Zhao also proved that $c_0(G)\le2|G|/5+1$ for
every group of rank two with $p\ge11$
\cite[Lemma~7]{[QLWZ2025]}.
For $G=C_p\oplus C_{np}$ with $n\ge2$, Qu and Han conjectured
that $c_0(G)=np+2p-3$ \cite[Concluding Remark]{[QH2017]}.

We also study the invariant $\mathsf N(G)$ introduced by Gao et al.
\cite{[GHHLYZ2022-1]}. A sequence $S$ over $G$ is weak-regular if
$\mathsf v_g(S)\le\ord(g)$ for every $g\in G$, where $\mathsf v_g(S)$
denotes the multiplicity of $g$ in $S$.
The invariant $\mathsf N(G)$ is the least integer $l$ such that every
weak-regular sequence $S$ over $G$ with $|S|\ge l$ has a nonempty
zero-sum subsequence $T$ satisfying $\mathsf v_g(T)=\mathsf v_g(S)$
for some $g\mid S$. Equivalently, $\supp(ST^{-1})\ne\supp(S)$.
Gao et al. determined
$\mathsf N(C_p^r)=r(p-1)+1$ for $p\le3$ and
$\mathsf N(C_p^2)=4p-2\lceil2\sqrt{p+1}\rceil+1$ for $p\ge3$
\cite[Theorem~1.3]{[GHHLYZ2022-2]}.
For cyclic groups, Jiang and the author proved
\[
2n-\lceil2\sqrt{n+1}\rceil
\le\mathsf N(C_n)
\le2n-\lceil2\sqrt{n+1}\rceil+1
\qquad(n\ge3),
\]
with each equality holding for infinitely many $n$
\cite[Theorem~1.1]{[JY2024]}.

Our bounds for additive bases follow from the estimate
\[
c_0(C_q\oplus C_{nq})\le
\max\left\{nq+2q^2-q-1,\,
\frac{nq(q+1)}p-q+p-2\right\},
\]
where $q$ is prime, $n\ge2$, and $p\ge3$ is the smallest prime
divisor of $nq$; see Theorem~\ref{thm:unified-c0}.
This estimate applies whether or not $q\mid n$.
The first two theorems give consequences for the two cases separately.

\begin{thm}\label{thm5.1}
	Let $q$ be prime and let $G=C_q\oplus C_{nq^2}$, where $n\ge1$ is not a power of $q$. If $p\ge3$ is the smallest prime dividing $|G|$, then
	\begin{align*}
		c_0(G)\le \max \{\frac{nq^2(q+1)}{p},nq^2 +f(q) -1\},
	\end{align*}
	where $f(q)= 4q(q-1) +1$ for $q\ge 7$, $f(q)= 2q^2 +3$ for $q\le 5$.
\end{thm}

\begin{thm}\label{thm1.2}
	Let $q$ be prime and let $G=C_q\oplus C_{nq}$, where $n\ge1$ and $q\nmid n$. If $p\ge3$ is the smallest prime dividing $|G|$, then
	
	\begin{align*}
		c_0(G)\le \max \{(2q-1)(n+q+1),\frac{nq^2}{p} +2\frac{nq}{p}\}. 
	\end{align*}
	
\end{thm}

For $\mathsf N(G)$, we prove the lower bound $r(n_r-1)+1$.
The other term in the following theorem is the direct-sum bound of
Gao et al. \cite[Theorem~1.1(2)]{[GHHLYZ2022-2]}.

\begin{thm}\label{thm1.4}
Let $G=C_{n_1}\oplus\cdots\oplus C_{n_r}$, where
$1<n_1\mid\cdots\mid n_r$. Then
\[
\mathsf{N}(G)\ge
\max\left\{r(n_r-1)+1,\ \sum_{i=1}^r\bigl(\mathsf{N}(C_{n_i})-1\bigr)+1\right\}.
\]
\end{thm}

\begin{thm}\label{thm1.5}
Let $G=C_p\oplus C_{np}$, where $n\ge2$ and $p\ge3$ is the
smallest prime dividing $|G|$. Then
\[
 2np-1\le\mathsf N(G)\le
 \begin{cases}
 2np-1,&p\mid n,\\
 5np-4,&p\ge5,\ p\nmid n,\\
 4np-3,&p=3,\ 3\nmid n.
 \end{cases}
\]
In particular, $\mathsf N(C_p\oplus C_{np^2})=2np^2-1$
whenever $n\ge1$ and $p\ge3$ is the smallest prime dividing the group order.
The two upper bounds for $p\nmid n$ also hold when $n=1$.
\end{thm}

Section~2 gives the notation and preliminary lemmas.
Section~3 proves Theorems~\ref{thm5.1} and~\ref{thm1.2}.
Section~4 proves Theorems~\ref{thm1.4} and~\ref{thm1.5}.
The final section discusses the remaining rank-two cases and two cyclic families.

\section{Notation and preliminaries}

Our notation for sequences follows the standard conventions of zero-sum
theory; see \cite{[GG2006],[GH2006]}.
A sequence over $G$ is finite and unordered, and repetitions are allowed.
For a sequence $S=g_1\cdots g_l$, its length is $|S|=l$, its sum is
$\sigma(S)=g_1+\cdots+g_l$, and its support is
$\supp(S)=\{g\in G:\mathsf v_g(S)>0\}$.
The relation $T\mid S$ means that $T$ is a subsequence of $S$;
then $ST^{-1}$ is the subsequence obtained by deleting $T$ from $S$.
The empty sequence is denoted by $1$ and has sum $0$.
Two subsequences $T_1,T_2$ of $S$ are disjoint if $T_1T_2\mid S$.
A homomorphism $\varphi:G\to G_1$ induces the sequence
$\varphi(S)=\varphi(g_1)\cdots\varphi(g_l)$ over $G_1$.

Let $\sum_0(S)=\{0\}\cup\sum(S)$.
Thus $\sum_0(S)$ includes the sum of the empty subsequence, whereas
$\sum(S)$ consists of sums of nonempty subsequences.
A nonempty sequence is a minimal zero-sum sequence if it has sum $0$
and no proper nonempty zero-sum subsequence.
The Davenport constant $\mathsf D(G)$ is the least integer $l$ such that
every sequence over $G$ of length at least $l$ has a nonempty zero-sum
subsequence. We use
$\mathsf D(C_{n_1}\oplus C_{n_2})=n_1+n_2-1$ for $1\le n_1\mid n_2$
\cite[Theorem~5.8.3]{[GH2006]}.
Let $\N$ denote the positive integers. We use the usual floor and
ceiling functions $\lfloor x\rfloor$ and $\lceil x\rceil$.

For each subset $A$ of $G$, denote by $\st(A) = \{g \in G : g + A = A\}$ the stabilizer of $A$ in $G$. Then $\st(A)$ is the maximal subgroup $H$ of $G$ such that $H + A = A$. We use the following form of Kneser's theorem; see \cite[Theorem~4.4]{[N1996]}.

\begin{lem}[Kneser]\label{lem2.1}
	Let $A_1,...,A_n$ be nonempty finite subsets of an
	abelian group $G$, and $H= \st(A_1+\cdots +A_n)$ be the stabilizer. Then
	$$|A_1+\cdots +A_n| \ge  |A_1+H|+\cdots +|A_n+H| -(n-1)|H|.$$
	
\end{lem}

\begin{lem}[Cauchy--Davenport]\label{lem2.2}
	Let $A_1,...,A_n$ be nonempty finite subsets of an
	abelian group $G=C_p$, where $p$ is prime. Then
	$$|A_1+\cdots +A_n|\ge\min\{p,\,|A_1|+\cdots +|A_n|-(n-1)\}.$$
	
\end{lem}

\begin{lem}\label{lem2.3}
	For every finite abelian group $G$,  
	
	$$m(G) \le c_0(G)\le |G|.$$
\end{lem}

See \cite[Lemma~2.3]{[GQZ2020]} for the lower bound and
\cite[Lemma~2.4]{[GHQQZ2015]} for the upper bound.

Qu and Han determined $c_0(G)$ for $G=C_p\oplus C_{p^n}$ as follows.

\begin{lem}[\cite{[QH2017]}]\label{lem2.4}
	Let $G=C_p\oplus C_{p^n}$, where $p$ is prime and $n\ge2$. Then
	$$c_0(G)=p^n+2p-3.$$
\end{lem}

\begin{lem}[{\cite[Lemma~2.3(1)]{[GHQQZ2015]}}]\label{lem5.2}
	Let $G$ be a finite abelian group, and let $p$ be the smallest prime dividing $|G|$. Let $S$ be a regular sequence over $G$ of length $|S|\ge \max \{|G|/p+p-2,~ \mathsf D(G)\}$. If $\sum (S)\ne G$, then $\st(\sum (S)) = \{0\}$.
\end{lem}

We use the cyclic bound
\[
\mathsf N(C_n)\le 2n-\lceil2\sqrt{n+1}\rceil+1\qquad(n\ge3)
\]
from \cite[Theorem~1.1]{[JY2024]}, and the value
$\mathsf N(C_p^2)=4p-2\lceil2\sqrt{p+1}\rceil+1$ for $p\ge3$
from \cite[Theorem~1.3(3)]{[GHHLYZ2022-2]}.
Also $\mathsf N(C_2)=2$.

\begin{lem}\label{lem2.7}
Let $G=C_p\oplus C_{np}$, where $p\ge3$ is the smallest prime divisor
of $|G|$. For $1<m\mid n$, let $H\le pG$ have order $m$, and let $k=n/m$.
Then
\[
2np-1-\bigl(\mathsf{N}(H)-1+(m-1)kp\bigr)\ge
\begin{cases}
2kp-1,&p\mid k,\ m\ne n,\\
4kp-3,&p=3,\ p\nmid k,\ m\ne n,\\
5kp-4,&p\ge5,\ p\nmid k,\ m\ne n,\\
\mathsf{N}(C_p^2),&m=n.
\end{cases}
\]
\end{lem}
\begin{proof}
Since $m\ge p\ge3$, the cyclic bound gives $\mathsf N(H)\le2m-3$.
Therefore
\[
\begin{aligned}
2np-1-\bigl(\mathsf N(H)-1+(m-1)kp\bigr)
&=kp(m+1)-\mathsf N(H)\\
&\ge kp(m+1)-2m+3.
\end{aligned}
\]
For $m<n$, subtracting $2kp-1$ from the last expression gives
$(kp-2)(m-1)+2\ge0$.
When $p=3$, subtracting $4kp-3$ gives $(m-3)(kp-2)\ge0$.
When $p\ge5$, subtracting $5kp-4$ gives
$(m-4)kp-2m+7\ge3m-13>0$.
These prove the first three estimates.
If $m=n$, then $k=1$, and the left side is at least
\[
(p-2)n+p+3\ge p^2-p+3>4p-7\ge\mathsf N(C_p^2).
\]
The strict inequality follows from
$p^2-5p+10=(p-2)(p-3)+4>0$.
\end{proof}

\begin{lem}\label{lem2.8}
Let $S$ be a weak-regular sequence over $G=C_n$, where $n\ge2$ and $|S|\ge n$.
Then either $\sum(S)=G$, or there is a zero-sum subsequence $T\mid S$
such that $\mathsf{v}_g(T)=\mathsf{v}_g(S)\ge1$ for some $g\mid S$.
\end{lem}
\begin{proof}
If $\mathsf{v}_g(S)=\ord(g)$ for some $g$, take $T=g^{\ord(g)}$.
Otherwise $0\notin\supp(S)$, and we may assume $\sum(S)\ne G$.
Since $|S|\ge \mathsf D(G)$, $\sum(S)=\sum_0(S)$.
Let $H=\st(\sum_0(S))$ and $m=|H|$.
If $m=1$, Kneser's theorem gives $|\sum_0(S)|\ge|S|+1>n$.
Thus $m>1$, and $\sum_0(S)\ne G$ gives $m<n$.
Since $n=dm$ for some integer $d\ge2$, Kneser's theorem also gives
\[
(|S|-|S_H|+1)m\le\bigl|\sum_0(S)\bigr|\le n-m,
\qquad
|S_H|\ge n-d+2\ge2m\ge\mathsf{N}(H).
\]
The middle inequality is equivalent to $(d-2)(m-1)\ge0$.
Apply the definition of $\mathsf{N}(H)$ to the weak-regular sequence $S_H$;
all occurrences in $S$ of the resulting element $g\in H$ belong to $S_H$.
\end{proof}

\begin{lem}\label{lem2.9}
Let $S,S'$ be sequences over a finite abelian group $G$, and let
$H\le G$ satisfy $H\subseteq\sum_0(S)$.
Suppose that $T_1\mid S'$ satisfies $\sigma(T_1)\in H$ and
$\mathsf{v}_g(T_1)=\mathsf{v}_g(S')\ge1$ for some
$g\in\supp(S')\setminus\supp(S)$.
Then $SS'$ has a zero-sum subsequence $T$ with
$\mathsf{v}_g(T)=\mathsf{v}_g(SS')$.
\end{lem}
\begin{proof}
Choose $T_2\mid S$ with $\sigma(T_2)=-\sigma(T_1)$ and take $T=T_1T_2$.
Since $g\notin\supp(S)$, its multiplicity in $T$ equals that in $SS'$.
\end{proof}

\section{Proof of Theorems \ref{thm5.1} and \ref{thm1.2}}

\begin{lem}\label{lem:c0-regular-sums}
Let $S$ be a regular sequence over a finite abelian group $G$. Then
\[
|\sum_0(S)|\ge \min\{|G|,|S|+1\}.
\]
In particular, $\sum_0(S)=G$ if $|S|\ge |G|-1$.
\end{lem}

\begin{proof}
Suppose that $|\sum_0(S)|<\min\{|G|,|S|+1\}$.
Let $H=\st(\sum_0(S))$ and $m=|H|$.
Then $H<G$ and $m\le |S|$, so Kneser's theorem gives
\[
|\sum_0(S)|\ge (|S|-|S_H|+1)m
\ge (|S|-m+2)m\ge |S|+1,
\]
a contradiction.
\end{proof}

\begin{lem}\label{lem:c0-short-block}
Let $K<H<G$, where $|K|\ge2$ and $H/K\cong C_q$ for a prime $q$.
Let $S$ be a sequence over $H\setminus K$ that is regular over $G$.
If $|S|\ge2q-1$, then there exists $T\mid S$ such that
$|T|\le q$ and $\sigma(T)\in K\setminus\{0\}$.
\end{lem}

\begin{proof}
Take $RR'\mid S$ with $|R|=q$ and $|R'|=q-1$.
Since $R$ is regular over $H$ and $|H|=q|K|\ge2q$,
Lemma~\ref{lem:c0-regular-sums} gives $|\sum_0(R)|\ge q+1$.
Thus there exist $R_1,R_2\mid R$ whose sums are distinct and belong to the same coset of $K$.
By Cauchy--Davenport, $\sum_0(R')+K=H$.
Choose $R_3\mid R'$ such that $\sigma(R_1R_3),\sigma(R_2R_3)\in K$.
Choose $U\in\{R_1R_3,R_2R_3\}$ with $\sigma(U)\ne0$.
Factor $U$ into subsequences whose images in $H/K$ are minimal zero-sum sequences.
Each factor has length at most $q$, and at least one has nonzero sum, as required.
\end{proof}

\begin{lem}\label{lem:c0-decomposition}
Let $K\le G$ with $|K|\ge2$ and $G/K\cong C_q^2$ for a prime $q$.
If $S$ is regular over $G$ and
\[
|S|\ge q|K|+2q^2-q-1,
\]
then $S=S_KT_1\cdots T_kR$, where $k=|K|-|S_K|$,
$|T_i|\le q$, $\sigma(T_i)\in K\setminus\{0\}$ and $\sum_0(R)+K=G$.
\end{lem}

\begin{proof}
Let $\varphi:G\to G/K$ be the natural homomorphism and let $k=|K|-|S_K|$.
Starting with $SS_K^{-1}$, choose at each step a subgroup of order $q$ containing the largest number of remaining images.
Before each of the first $k$ steps, the remaining length is at least
$|S|-|S_K|-q(k-1)\ge2q^2-1$, so the chosen subgroup contains at least $2q-1$ terms.
Apply Lemma~\ref{lem:c0-short-block} in its inverse image to obtain $T_1,\ldots,T_k$.
The remaining sequence $R=S(S_KT_1\cdots T_k)^{-1}$ has length
$|R|\ge2q^2-q-1$.
Suppose that $\sum_0(\varphi(R))\ne G/K$.
Since $|R|>(q+1)(q-2)$, some subgroup $\bar H$ of order $q$ contains at least $q-1$ terms of $\varphi(R)$.
For $H=\varphi^{-1}(\bar H)$, Cauchy--Davenport gives $\sum_0(\varphi(R_H))=\bar H$.
It also gives $|RR_H^{-1}|\le q-2$, since otherwise $\sum_0(\varphi(R))=G/K$.
Thus $|R_H|\ge2q^2-2q+1$.
Working backwards, suppose that $T_{j+1},\ldots,T_k$ lie in $H$.
Before $T_j$ was chosen, at most $q-2+q=2q-2$ terms lay outside $H$,
whereas at least $2q^2-2q+1$ lay in $H$.
Distinct subgroups of order $q$ in $G/K$ meet only at zero,
and the remaining terms have nonzero images. Thus the subgroup containing
the largest number of remaining images is $\bar H$, forcing $T_j$ to lie in $H$.
Hence all $T_i$ lie in $H$, and $|S_H|\ge|S|-q+2>|H|=q|K|$, a contradiction.
\end{proof}

\begin{lem}\label{lem:c0-KneserW0}
Let $p$ be the smallest prime dividing $|G|$, let $K<G$ with $|K|\ge2$, and let $l\ge2$ be an integer.
Suppose that $S$ is regular over $G$ and admits a decomposition
$S=S_KT_1\cdots T_kR$, where $k=|K|-|S_K|$,
$|T_i|\le l$, $\sigma(T_i)=t_i\in K\setminus\{0\}$ and $\sum_0(R)+K=G$.
If
\[
|S|\ge\frac{|G|}{p}+l\frac{|K|}{p}-l+p-2,
\]
then $\sum(S)=G$.
\end{lem}

\begin{proof}
Let $W=S_Kt_1\cdots t_k$. Since $|W|=|K|\ge \mathsf D(K)$, we have $0\in\sum(S)$.
If $W$ is regular over $K$, then $\sum_0(W)=K$ and the conclusion follows.
Otherwise choose a minimal subgroup $H_1<K$ with $|W_{H_1}|\ge|H_1|=m$.
Then $H_1\ne\{0\}$ and $W_{H_1}$ is regular over $H_1$.
Let $r=m-|S_{H_1}|-1$.
After reordering, $S_{H_1}t_1\cdots t_r$ is regular over $H_1$ of length $m-1$.
Its subsequence sums together with $0$ therefore cover $H_1$ by Lemma~\ref{lem:c0-regular-sums}.
Let $W_0=St_1\cdots t_r(T_1\cdots T_r)^{-1}$ and $H_2=\st(\sum_0(W_0))$.
Since the $T_i$ are disjoint, replacing each selected occurrence of
$t_i$ by $T_i$ shows that $\sum_0(W_0)\subseteq\sum_0(S)$.
Together with $0\in\sum(S)$, this gives $\sum_0(W_0)\subseteq\sum(S)$.
The subsequence $S_{H_1}t_1\cdots t_r$ of $W_0$ has
$\sum_0(S_{H_1}t_1\cdots t_r)=H_1$, so $\sum_0(W_0)$ is a union
of cosets of $H_1$. Hence $H_1\le H_2$.
We may assume $H_2<G$; hence $p\le|H_2|\le|G|/p$.
Since $r\le|K|/p-1$, $|W_0|\ge|S|-(l-1) r$ and
$|(W_0)_{H_2}|\le|S_{H_2}|+r\le|H_2|-1+r$, Kneser's theorem gives
\[
|\sum_0(W_0)|\ge(|S|-lr+2-|H_2|)|H_2|
\ge(|G|/p+p-|H_2|)|H_2|\ge|G|.
\]
The last inequality follows from $(|H_2|-p)(|G|/p-|H_2|)\ge0$, and thus $\sum(S)=G$.
\end{proof}

\begin{thm}\label{thm:unified-c0}
Let $G=C_q\oplus C_{nq}$, where $q$ is prime and $n\ge2$,
and let $p\ge3$ be the smallest prime dividing $|G|$. Then
\[
c_0(G)\le
\max\left\{nq+2q^2-q-1,\,
\frac{nq(q+1)}p-q+p-2\right\}.
\]
\end{thm}

\begin{proof}
Let $S$ be regular over $G$ with length at least the stated bound.
For $K=qG$, we have $|K|=n$ and $G/K\cong C_q^2$.
The first term of the bound allows us to apply Lemma~\ref{lem:c0-decomposition}.
The second term is the bound in Lemma~\ref{lem:c0-KneserW0} with $l=q$.
Hence $\sum(S)=G$.
\end{proof}

\begin{proof}[Proof of Theorem~\ref{thm5.1}]
Apply Theorem~\ref{thm:unified-c0} with $nq$ in place of $n$. It gives
\[
c_0(G)\le\max\left\{nq^2+2q^2-q-1,\,
\frac{nq^2(q+1)}p-q+p-2\right\}.
\]
Since $p\le q$, the second term is at most $nq^2(q+1)/p$.
For the first, $f(q)-2q^2+q$ equals $2q^2-3q+1$ when $q\ge7$
and $q+3$ when $q=3,5$, and is positive in both cases.
This proves the stated bound.
\end{proof}

\begin{proof}[Proof of Theorem~\ref{thm1.2}]
If $n=1$, then $G=C_q^2$ and $c_0(G)=2q-1$, which is at most
$(2q-1)(n+q+1)$.
Suppose that $n\ge2$. Since $p\ge3$ is the smallest prime dividing
$|G|$, Theorem~\ref{thm:unified-c0} applies and gives
\[
c_0(G)\le\max\left\{nq+2q^2-q-1,\,
\frac{nq(q+1)}p-q+p-2\right\}.
\]
The first is at most $(2q-1)(n+q+1)$, since subtracting
$nq+2q^2-q-1$ from that expression gives $(q-1)n+2q>0$.
For the second, $p\le q$ implies
$nq(q+1)/p-q+p-2\le nq(q+1)/p$.
The latter is at most $nq^2/p+2nq/p$, because their difference is
$nq/p>0$.
Taking the maximum of these two comparisons proves the result.
\end{proof}

\section{Proof of Theorems \ref{thm1.4} and \ref{thm1.5}}

\begin{proof}[Proof of Theorem~\ref{thm1.4}]
Choose a basis $(x_1,\ldots,x_r)$ with $\ord(x_i)=n_i$ and let
$g_0=x_r$, $g_i=x_r+x_i$ for $1\le i<r$.
All $g_i$ have order $n_r$, so $S=\prod_{i=0}^{r-1}g_i^{n_r-1}$ is weak-regular.
For a zero-sum subsequence $T=\prod_{i=0}^{r-1}g_i^{a_i}$, the coordinates give
$a_i\equiv0\pmod{n_i}$ for $i\ge1$ and $\sum_i a_i\equiv0\pmod{n_r}$.
Thus also $a_0\equiv0\pmod{n_1}$, and none of the $a_i$ equals $n_r-1$.
This proves $\mathsf{N}(G)\ge r(n_r-1)+1$.
For completeness, consider the direct-sum bound from
\cite[Theorem~1.1(2)]{[GHHLYZ2022-2]}.
For each $i$, choose a weak-regular sequence over $\langle x_i\rangle$
of length $\mathsf N(C_{n_i})-1$ having no nonempty zero-sum subsequence
that contains all occurrences of any of its elements.
None contains $0$, and the direct factors intersect only at $0$, so their
product in $G$ is weak-regular.
In a direct sum, a sequence has sum zero if and only if its sum in every
coordinate is zero. A zero-sum subsequence containing all occurrences of
some element would therefore have the same property in the corresponding
factor, contrary to the choice of the sequences. The product has length
$\sum_{i=1}^r(\mathsf N(C_{n_i})-1)$, giving the second bound.
\end{proof}

If $rn_r\ge2\sum_{i=1}^r n_i$, the first lower bound dominates the second,
since $\mathsf{N}(C_{n_i})-1\le2n_i-2$.

\begin{lem}\label{lem6.1}
Let $G=C_p\oplus C_{np}$, where $n\ge2$ and $p\ge3$ is the
smallest prime dividing $|G|$. Then
\[
c_0(G)\le
\begin{cases}
\left\lceil(4np-2)/3\right\rceil,&p\mid n,\\
\left\lceil(25np-20)/6\right\rceil,&p\ge5,\ p\nmid n,\\
\left\lceil(16np-12)/5\right\rceil,&p=3,\ 3\nmid n.
\end{cases}
\]
\end{lem}

\begin{proof}
Suppose first that $p\mid n$.
If $n$ is a power of $p$, then
$c_0(G)=np+2p-3\le\lceil(4np-2)/3\rceil$, since
$np\ge p^2\ge6p-9$.
Otherwise $n\ge p(p+2)$, whence
$np\ge p^2(p+2)\ge6p^2-3p-1$. Thus
\[
np+2p^2-p-1\le\frac{4np-2}{3},
\qquad
n(p+1)-2\le\frac{4np-2}{3}.
\]
Theorem~\ref{thm:unified-c0} proves the first assertion.
If $p\nmid n$, then $n\ge p+2$, and the same theorem gives
$c_0(G)\le3np$.
For $p\ge5$, this is at most $\lceil(25np-20)/6\rceil$;
for $p=3$, it is at most $\lceil(16np-12)/5\rceil$,
since $np\ge15$.
\end{proof}

\begin{prop}\label{prop:regular-N-reduction}
Let $G=C_p\oplus C_{np}$, where $p$ is prime and $n\ge1$, and let $l_0$ be a positive integer.
Let $r=\lceil l_0/(np-1)\rceil$. If
\[
l_0-\lfloor l_0/r\rfloor\ge c_0(G),
\]
then every regular sequence $S$ over $G$ with $|S|\ge l_0$
has a nonempty zero-sum subsequence $T$ satisfying
$\mathsf v_g(T)=\mathsf v_g(S)$ for some $g\mid S$.
\end{prop}

\begin{proof}
Since $\langle g\rangle<G$ for every $g\in G$,
regularity gives $\mathsf v_g(S)\le\ord(g)-1\le np-1$.
Hence $|\supp(S)|\ge r$.
Choose $g_0\mid S$ with minimum multiplicity
$a=\mathsf v_{g_0}(S)$. The sequence $S'=Sg_0^{-a}$
contains none of the occurrences of $g_0$.
Then $S'$ is regular, $a\le\lfloor |S|/r\rfloor$, and
\[
|S'|\ge |S|-\lfloor |S|/r\rfloor
\ge l_0-\lfloor l_0/r\rfloor\ge c_0(G).
\]
Choose $T'\mid S'$ with $\sigma(T')=-ag_0$.
Then $T=T'g_0^a$ has the required property.
\end{proof}

For $n\ge2$ and $p\ge3$ the smallest prime dividing $|G|$, let
\[
\operatorname{ub}(G)=\begin{cases}
2np-1,&p\mid n,\\
5np-4,&p\ge5,\ p\nmid n,\\
4np-3,&p=3,\ 3\nmid n.
\end{cases}
\]

\begin{cor}\label{cor:regular-main}
Let $G=C_p\oplus C_{np}$, where $n\ge2$ and $p\ge3$ is the
smallest prime dividing $|G|$. Every regular sequence $S$ over $G$
with $|S|\ge \operatorname{ub}(G)$ has a nonempty zero-sum subsequence $T$ such that
$\mathsf v_g(T)=\mathsf v_g(S)$ for some $g\mid S$.
\end{cor}
\begin{proof}
Apply Proposition~\ref{prop:regular-N-reduction} with $l_0=\operatorname{ub}(G)$.
The corresponding values of $r$ are $3$, $6$, and $5$,
according as $p\mid n$, $p\ge5$ and $p\nmid n$, or
$p=3$ and $3\nmid n$.
In each case, the required inequality is precisely
Lemma~\ref{lem6.1}.
\end{proof}

\begin{lem}\label{lem:4.0}
Let $S$ be a non-regular sequence over a finite abelian group $G$.
Then there exists a proper subgroup $H<G$ such that $S_H$
is regular over $H$, $|S_H|\ge|H|$, and $\sum(S_H)=H$.
\end{lem}

\begin{proof}
Choose $H<G$ of minimum order subject to $|S_H|\ge|H|$.
For every proper subgroup $H'<H$, minimality gives
$|S_{H'}|\le|H'|-1$. Thus $S_H$ is regular over $H$.
Since $c_0(H)\le|H|$, we obtain $\sum(S_H)=H$.
\end{proof}

\begin{lem}\label{lem:N-heavy-fiber}
Let $S=S_HS_1$ with $\sum_0(S_H)=H$ and
$\supp(S_1)\subseteq G\setminus H$, and let $\Phi:G\to G/H$
be the natural homomorphism.
Suppose that some $\bar g\in\supp(\Phi(S_1))$ satisfies
$\mathsf v_{\bar g}(\Phi(S_1))\ge\ord(\bar g)$ and that
$\Phi^{-1}(\bar g)\cap\supp(S_1)$ contains at least two elements.
Then $S$ has a nonempty zero-sum subsequence $T$ with
$\mathsf v_g(T)=\mathsf v_g(S)$ for some $g\mid S_1$.
\end{lem}

\begin{proof}
Let $V\mid S_1$ consist of all terms whose image is $\bar g$,
and let $n=\ord(\bar g)$.
If $\mathsf v_h(V)\ge n$ for some $h\mid V$, let $|V|=kn+t$,
where $0\le t<n$, and take $T_1=Vh^{-t}$.
Then $\sigma(T_1)\in H$, and
$\mathsf v_g(T_1)=\mathsf v_g(S_1)$ for every $g\mid V$ with $g\ne h$.
Such an element exists by hypothesis.
Otherwise, choose $g\mid V$. Since $\mathsf v_g(V)<n\le|V|$,
there is an $n$-term subsequence $T_1\mid V$ containing all occurrences
of $g$. Again $\sigma(T_1)\in H$ and
$\mathsf v_g(T_1)=\mathsf v_g(S_1)$.
Lemma~\ref{lem2.9} completes the proof in both cases.
\end{proof}

\begin{lem}\label{lem:N-index-p-generators}
Let $G=C_p\oplus C_{np}$, where $p\ge3$ is the smallest prime
dividing $|G|$, $n\ge5$, and $H<G$ is cyclic of order $np$.
Let $S=S_Hg_0^b$ satisfy $g_0\notin H$, $\ord(g_0)=np$,
$b\le np-1$, $|S|\ge2np-1$, and $|S_H|\ge np$.
Suppose that every term of $S_H$ generates $H$.
Choose $h\mid S_H$ with minimum multiplicity $a$.
Let $R=S_Hh^{-a}$. If $|R|\le np-1$ and $a<np/2$,
then $S$ has a nonempty zero-sum subsequence $T$ with $\mathsf v_h(T)=\mathsf v_h(S)$.
\end{lem}

\begin{proof}
Let $b=ps_g+t_g$, and
$R=\prod_{i=1}^{\ell-1}g_i^{ps_i+t_i}$, where $0\le t_i,t_g<p$.
First suppose that $a\ge2p$.
As $\ell\le1+(np-1)/a$, convexity of
$x+(p-1)(np-1)/x$ on $[2p,np/2]$ gives
\[
p\Big(s_g+\sum_i s_i\Big)
\ge2np-1-a-\ell(p-1)
\ge3np/2-3p\ge np-p/2.
\]
Thus $s_g+\sum_i s_i\ge n$.
Since $s_g\le n-1$, choose $i$ with $s_i>0$.
Consider $U=(pg_0)^{s_g}\prod_j(pg_j)^{s_j}$ and delete one
occurrence of $pg_i$.
The remaining sequence has at least $n-1$ terms, all generating
$pH=pG$, so its $\sum_0$ equals $pH$ by
Lemma~\ref{lem:c0-regular-sums}.
The deleted occurrence of $pg_i$ leaves $g_i^p$ unused in the corresponding
subsequences of $Rg_0^{ps_g}$. Since $\sum_0(g_i^p)+pH=H$, we obtain
$H\subseteq\sum_0(Rg_0^{ps_g})$.

Now suppose that $a\le2p-1$.
Then $s_g\ge n-2$ and $|R|\ge(n-2)p+1\ge3p+1$.
Take disjoint $U,V\mid R$ with $|U|=|V|=p$.
Kneser's theorem gives $|\sum_0(U)|\ge p+1$.
Hence two different sums from $U$ have the same image in $H/pH$.
Using $\sum_0(V)+pH=H$, extend both by the same subsequence of $V$
to sums in $pH$. At least one is nonzero.
Thus there is $W\mid UV$ with $\sigma(W)\in pH\setminus\{0\}$.
Since $pg_0$ generates $pH$ and $\sigma(W)\ne0$, every proper
subgroup of $pH$ contains at most one term of $(pg_0)^{s_g}\sigma(W)$,
and the trivial subgroup contains none. This sequence is therefore regular
over $pH$ and has length at least $n-1$, so its $\sum_0$ equals $pH$.
Also $|RW^{-1}|\ge p+1$, and its image has all terms nonzero in $H/pH$.
Again $H\subseteq\sum_0(Rg_0^{ps_g})$.
In either case choose $T_1\mid Rg_0^{ps_g}$ with
$\sigma(T_1)=-ah$ and take $T=T_1h^a$.
\end{proof}

\begin{lem}\label{lem:N-index-p-low-order}
Let $G=C_p\oplus C_{np}$, where $n\ge2$ and $p\ge3$ is the
smallest prime dividing $|G|$, and let $H<G$ be cyclic of order $np$.
Let $S=S_Hg_0^b$ satisfy $g_0\notin H$, $\ord(g_0)=np$,
$b\le np-1$, $|S|\ge2np-1$, and suppose that $S_H$ is regular over $H$.
If $\ord(h')\le n$ for some $h'\mid S_H$, then $S$ has a nonempty zero-sum
subsequence $T$ with $\mathsf v_{h'}(T)=\mathsf v_{h'}(S)$.
\end{lem}

\begin{proof}
Let $a=\mathsf v_{h'}(S)$, $R'=S_H(h')^{-a}$,
$s=\lfloor b/p\rfloor$, and $U=R'(pg_0)^s$.
Regularity gives $a\le n-1$, while $s\le n-1$ and
$b-s\le n(p-1)$. Hence $|U|=|S|-a-b+s\ge np$.
Suppose that $\sum_0(U)\ne H$, and let
$H_1=\st(\sum_0(U))$ and $d=|H_1|$.
Since $0\nmid U$, Kneser's theorem implies $H_1\ne\{0\}$.
It also gives
\[
|U|-|U_{H_1}|\le np/d-2,\qquad
|U_{H_1}|\le d-1+s\le d+n-2.
\]
As $p\le d\le n$,
$|U|\le np/d+d+n-4\le3n-4<np$, a contradiction.
Thus $\sum_0(U)=H\subseteq\sum_0(R'g_0^{ps})$.
Choose $T_1\mid R'g_0^{ps}$ with $\sigma(T_1)=-ah'$;
then $T=T_1(h')^a$ has the required property.
\end{proof}

\begin{lem}\label{lem:N-noncyclic-subgroup-three}
Let $H=C_3\oplus C_{3m}$, where $m\ge1$ is odd.
Every regular sequence $R$ over $H$ of length at least $|H|=9m$
has a nonempty zero-sum subsequence $T$ with
$\mathsf v_g(T)=\mathsf v_g(R)$ for some $g\mid R$.
\end{lem}

\begin{proof}
Every multiplicity is at most $3m-1$, so $|\supp(R)|\ge4$.
Choose $g\mid R$ of minimum multiplicity $a$ and let $R'=Rg^{-a}$.
Then $R'$ is regular over $H$ and
$|R'|\ge\lceil3|R|/4\rceil\ge\lceil27m/4\rceil$.
For $m\ge5$, Theorem~\ref{thm:unified-c0} gives
\[
c_0(H)\le\max\{3m+14,4m-2\}\le6m\le\lceil27m/4\rceil.
\]
For $m=1,3$, use $c_0(C_3^2)=5$ and $c_0(C_3\oplus C_9)=12$.
Thus $\sum(R')=H$. Choose $T'\mid R'$ with $\sigma(T')=-ag$.
Then $T=T'g^a$ is zero-sum and $\mathsf v_g(T)=\mathsf v_g(R)$.
\end{proof}

\begin{lem}\label{lem:N-half-quotient}
Let $H=C_m$, let $R$ be regular over $H$ with $|R|\ge m/2$,
and let $H_0\le H$ be nontrivial. Then $\sum_0(R)+H_0=H$.
\end{lem}
\begin{proof}
Suppose otherwise and let $H_1$ be the inverse image in $H$ of the
stabilizer of the image of $\sum_0(R)$ in $H/H_0$.
For $k=|H_1|$, we have $2\le k\le m/2$.
Kneser's theorem and regularity give
\[
 |R|-|R_{H_1}|\le m/k-2,\qquad
 |R|\le k+m/k-3\le m/2-1,
\]
contrary to $|R|\ge m/2$. The last inequality follows by convexity
of $x+m/x$ on $[2,m/2]$.
\end{proof}

\begin{lem}\label{lem:N-kernel-blocks}
Let $S=S_HS_1$ be a sequence over an abelian group $G$, where
$H=C_m$, $m\ge2$, $S_H$ is regular over $H$, $|S_H|\ge m$,
and every multiplicity in $S_H$ is smaller than the corresponding element order.
Suppose that pairwise disjoint nonempty subsequences of $S_1$ have sums
forming a sequence $U$ over $H$.
If $\sum_0(U)$ contains a nontrivial subgroup $H_0\le H$, then $S$
has a nonempty zero-sum subsequence $T$ with
$\mathsf v_h(T)=\mathsf v_h(S)$ for some $h\mid S_H$.
In particular, this holds if $U$ has no zero term, is supported in a
nontrivial subgroup $K\le H$, and $|U|\ge|K|-1$.
\end{lem}
\begin{proof}
The strict multiplicity bound and $|S_H|\ge m$ ensure that $S_H$
has at least two distinct elements. Choose $h\mid S_H$ of minimum multiplicity $a$. Then $R=S_Hh^{-a}$ has length $|R|\ge m/2$.
Lemma~\ref{lem:N-half-quotient} gives $\sum_0(R)+H_0=H$.
Choose $V\mid R$ with $\sigma(V)+ah\in H_0$.
Since $H_0\subseteq\sum_0(U)$, a subsequence of $U$ has sum
$-\sigma(V)-ah$. The corresponding pairwise disjoint subsequences of
$S_1$ form a subsequence $T_1$ with the same sum.
Thus $T=Vh^aT_1$ is zero-sum. Since $h\nmid R$ and $h\nmid S_1$,
we also have $\mathsf v_h(T)=a=\mathsf v_h(S)$.
For the last assertion, a nontrivial stabilizer of $\sum_0(U)$
is itself contained in $\sum_0(U)$.
If the stabilizer is trivial, Kneser's theorem gives
$|\sum_0(U)|\ge|U|+1\ge|K|$, so $\sum_0(U)=K$.
\end{proof}

\begin{lem}\label{lem:N-cyclic-quotient}
Let $G=C_p\oplus C_{np}$, where $n\ge2$ and $p\ge3$ is the
smallest prime dividing $|G|$.
Let $S=S_HS_1$ satisfy $|S|\ge2np-1$ and
$\mathsf v_g(S)<\ord(g)$ for every $g\mid S$.
Suppose that $H<G$ is cyclic of order $m$, that $S_H$ is regular over $H$,
and that $m\le|S_H|\le\mathsf N(H)-1$.
If $G/H$ is cyclic of order greater than $p$, then $S$ has a nonempty
zero-sum subsequence $T$ with $\mathsf v_g(T)=\mathsf v_g(S)$
for some $g\mid S$.
\end{lem}
\begin{proof}
By Lemma~\ref{lem:c0-regular-sums}, $\sum_0(S_H)=H$.
Let $k=np/m>1$, so that $|G/H|=pk$, and let $\Phi:G\to G/H$
be the natural homomorphism.
Since $G/pG\cong C_p^2$ and $(G/H)/p(G/H)\cong C_p$,
$H/(H\cap pG)$ has order $p$. Thus $p\mid m$ and,
as $H$ is cyclic, $H\cap pG=pH$.
We have $k\ge p$ and $|S_1|\ge2np-\mathsf N(C_m)$.

By Lemma~\ref{lem:N-heavy-fiber}, we may assume that
$\Phi^{-1}(\bar g)\cap\supp(S_1)$ contains exactly one element whenever
$\mathsf v_{\bar g}(\Phi(S_1))\ge\ord(\bar g)$.
For each $g\mid S_1$, let $r_g=\ord(\Phi(g))$ and choose
$\lfloor\mathsf v_g(S_1)/r_g\rfloor$ disjoint copies of $g^{r_g}$.
Let $T_0$ be their product, $U$ the sequence of their sums, and
$R=S_1T_0^{-1}$.
Every term of $U$ is nonzero, since a selected subsequence $g^{r_g}$
satisfies $r_g\le\mathsf v_g(S)<\ord(g)$.
The choice of $T_0$ also ensures that $\Phi(R)$ is weak-regular.

Suppose that $R_1\mid R$ satisfies $\sigma(\Phi(R_1))=0$ and
$\mathsf v_{\bar g}(\Phi(R_1))=\mathsf v_{\bar g}(\Phi(R))$
for some $\bar g\mid\Phi(R)$.
Then $R_1$ contains every term of $R$ whose image is $\bar g$.
Consequently $T_1=R_1T_0$ has sum in $H$ and
$\mathsf v_h(T_1)=\mathsf v_h(S_1)$ for every
$h\in\Phi^{-1}(\bar g)\cap\supp(S_1)$.
Lemma~\ref{lem2.9} gives the required subsequence of $S$.
In particular, it suffices to prove $|R|\ge\mathsf N(G/H)$.

Suppose first that $m>p$, so $m\ge p^2$.
By Lemma~\ref{lem:N-kernel-blocks}, we may assume that $\sum_0(U)$
contains no nontrivial subgroup of $H$.
Let $U_1=U_{pH}$ and $U_2=UU_1^{-1}$, with $a=|U_1|$ and $b=|U_2|$.
The same lemma, applied to $U_1$ with $K=pH$, gives $a\le m/p-2$.
Each term of $U_1$ is the sum of a selected subsequence of length at most $pk$.
If a selected subsequence corresponding to a term of $U_2$ had length
divisible by $p$, its sum would lie in $H\cap pG=pH$.
Its length therefore divides $k$ and is not divisible by $p$, so is at most $k$.
Since its length is greater than one, $U_2\ne1$ implies $k\ge p+2$.

If $b\le p-2$, then
\[
 |T_0|\le k(pa+b)\le np-(p+2)k,
 \qquad |R|\ge m(k-2)+(p+2)k\ge2pk.
\]
Here $\mathsf N(C_m)\le2m$, and
$m(k-2)\ge k(p-2)$ follows from $m\ge p$ and $k\ge p$.
Thus $|R|\ge2pk\ge\mathsf N(C_{pk})$.

If $b\ge p-1$, choose $V\mid U_2$ of length $p-1$.
Its image in $H/pH=C_p$ has all terms nonzero, so
$\sum_0(V)+pH=H$.
The stabilizer of $\sum_0(U)$ is trivial, since otherwise it would be a
nontrivial subgroup contained in $\sum_0(U)$.
For every $W\mid U$, the identity
$\sum_0(U)=\sum_0(W)+\sum_0(UW^{-1})$ gives
$\st(\sum_0(W))\le\st(\sum_0(U))=\{0\}$.
Kneser's theorem therefore yields $|\sum_0(U_1)|\ge a+1$.
Since $\sum_0(V)$ meets every coset of $pH$,
$|\sum_0(U_1V)|\ge p(a+1)$.
Each remaining term of $U_2$ increases the number of subsequence sums
by at least one. Otherwise that nonzero term would belong to the
stabilizer of the preceding sumset, which is trivial.
It follows that
\[
 pa+b+1\le|\sum_0(U)|\le m-1,\qquad |T_0|\le k(pa+b)\le np-2k.
\]
Hence $|R|\ge m(k-2)+2k\ge2pk$, since $m\ge p^2$, $k\ge p+2$, and
$p^2(k-2)-2k(p-1)\ge(p-2)(p^2-2)>0$.
Again $|R|\ge\mathsf N(C_{pk})$.

It remains to consider $m=p$.
Since $pH=\{0\}$ and every term of $U$ is nonzero, every selected
subsequence has length at most $k$.
If $|U|\ge p-2$, then $|\sum_0(U)|\ge p-1$ by Cauchy--Davenport.
Choose distinct $h,x\in\supp(S_H)$; such a choice exists because
$|S_H|\ge p$ and every multiplicity is less than $p$.
Another application of Cauchy--Davenport gives
$\sum_0(U)+\{0,x\}=H$.
Since $\sum_0(U)\subseteq\sum_0(T_0)$, there is $T_1\mid T_0x$
with $\sigma(T_1)=-\mathsf v_h(S)h$.
Then $T=T_1h^{\mathsf v_h(S)}$ is zero-sum and
$\mathsf v_h(T)=\mathsf v_h(S)$.
We may therefore assume $|U|\le p-3$.

If $T_0=1$, then $R=S_1$ and
\[
 |R|\ge2pk-\mathsf N(C_p)\ge2pk-2p\ge\mathsf N(C_{pk}).
\]
The last inequality uses $pk\ge p^2$ and the cyclic bound
$\mathsf N(C_{pk})\le2pk-\lceil2\sqrt{pk+1}\rceil+1$.
Otherwise choose $g\mid S_1$ with $\mathsf v_g(S_1)\ge r_g$.
Let $t=\mathsf v_g(R)<r_g\le k$ and $R'=Rg^{-t}$.
By the assumption on the cosets, $\Phi(R')$ contains no occurrence of
$\Phi(g)$. It is weak-regular and
\[
 |R'|\ge2pk-\mathsf N(C_p)-(p-3)k-(k-1)
 =(p+2)k-\mathsf N(C_p)+1\ge pk.
\]
Apply Lemma~\ref{lem2.8} to $\Phi(R')$.
If $\sum(\Phi(R'))=G/H$, choose $R_1\mid R'$ such that
$\Phi(\sigma(R_1))=-\mathsf v_g(S_1)\Phi(g)$.
Then $T_1=R_1g^{\mathsf v_g(S_1)}$ has sum in $H$ and
$\mathsf v_g(T_1)=\mathsf v_g(S_1)$, so Lemma~\ref{lem2.9} applies.
Otherwise there is $R_1\mid R'$ with $\sigma(\Phi(R_1))=0$ and
$\mathsf v_{\bar h}(\Phi(R_1))=\mathsf v_{\bar h}(\Phi(R'))$
for some $\bar h\mid\Phi(R')$.
Since $\bar h\ne\Phi(g)$, its multiplicity in $\Phi(R')$ equals that
in $\Phi(R)$. The argument for $R_1T_0$ above completes the proof.
\end{proof}

\begin{lem}\label{lem:N-index-p-complete}
Let $G=C_p\oplus C_{np}$, where $n\ge2$ and $p\ge3$ is the
smallest prime dividing $|G|$.
Let $S$ satisfy $|S|\ge2np-1$ and
$\mathsf v_g(S)<\ord(g)$ for every $g\in\supp(S)$.
Suppose that $H<G$ is cyclic of order $np$, that $S_H$ is regular
over $H$, and that $np\le|S_H|\le\mathsf N(H)-1$.
Then $S$ has a nonempty zero-sum subsequence $T$ such that
$\mathsf v_g(T)=\mathsf v_g(S)$ for some $g\in\supp(S)$.
\end{lem}

\begin{proof}
Let $S_1=SS_H^{-1}$, and let $\Phi:G\to G/H\cong C_p$ be the quotient map.
Lemma~\ref{lem:c0-regular-sums} gives $\sum_0(S_H)=H$.
By Lemma~\ref{lem2.9}, it suffices to find $T_1\mid S_1$
with $\sigma(T_1)\in H$ and $\mathsf v_g(T_1)=\mathsf v_g(S_1)$
for some $g\mid S_1$.
Since $n\ge p$ and $np\ge p^2$, the cyclic upper bound gives
\[
 |S_1|\ge2np-\mathsf N(C_{np})
 \ge\lceil2\sqrt{np+1}\rceil-1\ge2p\ge\mathsf N(C_p).
\]
If $\Phi(S_1)$ is weak-regular, its length gives a subsequence
$T_1\mid S_1$ such that $\sigma(\Phi(T_1))=0$ and
$\mathsf v_{\bar g}(\Phi(T_1))=\mathsf v_{\bar g}(\Phi(S_1))$
for some $\bar g\mid\Phi(S_1)$.
Thus $\sigma(T_1)\in H$, and $T_1$ contains all occurrences of every
element of $S_1$ whose image is $\bar g$. Lemma~\ref{lem2.9} applies.

Otherwise some $\bar g\in G/H$ has multiplicity greater than $p$
in $\Phi(S_1)$. If $\Phi^{-1}(\bar g)\cap\supp(S_1)$ contains at least
two elements, apply Lemma~\ref{lem:N-heavy-fiber}.
Suppose that this intersection is $\{g\}$.
If $S_1$ has another element $g_0\ne g$, then
$\sum_0(g^pS_H)=G$, because $g+H$ generates $G/H$ and
$\sum_0(S_H)=H$.
Choose $T'\mid g^pS_H$ with $\sigma(T')=-\mathsf v_{g_0}(S_1)g_0$.
Since $g_0\nmid g^pS_H$, the sequence
$T=T'g_0^{\mathsf v_{g_0}(S_1)}$ is a zero-sum subsequence of $S$
with $\mathsf v_{g_0}(T)=\mathsf v_{g_0}(S)$.

It remains to consider $S_1=g_0^b$, where $b\ge p$ and
$b<\ord(g_0)\le np$.
Choose $h\mid S_H$ of minimum multiplicity $a$, and let $R=S_Hh^{-a}$.
The strict multiplicity bounds and $|S_H|\ge np$ ensure that
$S_H$ has at least two distinct elements, so $|R|\ge|S_H|/2$.
If $|R|\ge np$, Lemma~\ref{lem:c0-regular-sums} gives $\sum_0(R)=H$.
Choose $T_1\mid R$ with $\sigma(T_1)=-ah$; then $T=T_1h^a$
has $\mathsf v_h(T)=a=\mathsf v_h(S)$.
Hence assume $|R|\le np-1$.

\smallskip\noindent\emph{Case 1: $|\supp(R)|\ge2$.}
Then $2a\le|R|\le np-1$ and
$b=|S|-|R|-a\ge np-a>np/2$, forcing $\ord(g_0)=np$.
If some $h'\mid S_H$ does not generate $H$, then
$\ord(h')\le np/p=n$, so Lemma~\ref{lem:N-index-p-low-order} applies.
We may therefore assume that every term of $S_H$ generates $H$.
For $n\ge5$, apply Lemma~\ref{lem:N-index-p-generators}.
The only remaining possibility is $p=n=3$.
Here $Rg_0^b$ is regular over $G$: a proper subgroup containing
any term of $R$ must equal $H$, where $|R|\le8$; any other proper
subgroup contains no terms or only $b$ occurrences of $g_0$,
and in the latter case has order at least $\ord(g_0)>b$.
Since $a\le4$, we have $|Rg_0^b|\ge13\ge c_0(C_3\oplus C_9)=12$.
Hence there is $T_1\mid Rg_0^b$ with $\sigma(T_1)=-ah$,
and $T=T_1h^a$ has the required property.

\smallskip\noindent\emph{Case 2: $|\supp(R)|=1$.}
Then $R=f^c$ for some $f\in H$, with $a\le c\le np-1$.
Since $2c\ge|S_H|\ge np$ and $c<\ord(f)$, we have $\ord(f)=np$.
Let $c=ps_f+t_f$ and $b=ps_g+t_g$, where $0\le t_f,t_g<p$.
If $a\ge(n-1)p$, then $s_f\ge n-1$ and $s_g\ge1$.
Otherwise $a\le(n-1)p-1$, so $c+b\ge(n+1)p$ and
\[
 s_f+s_g\ge
 \left\lceil\frac{(n+1)p-2(p-1)}p\right\rceil=n.
\]
Thus $s_f+s_g\ge n$ in both cases, and $s_f\ge\lfloor n/2\rfloor\ge1$.
Consider the sequence
\[
 V=(pf)^{s_f-1}(pg_0)^{s_g}\quad\text{over }pH,
 \qquad |V|\ge n-1.
\]
Indeed $pH=pG\cong C_n$ and $pf$ generates this group.
Since $g_0+H$ has order $p$, we have $p\mid\ord(g_0)$, and the integer
$d=\ord(pg_0)=\ord(g_0)/p$ satisfies $d>1$ and $s_g\le d-1$.
No proper subgroup of $pH$ contains $pf$; one containing $pg_0$
has order at least $d$ and contains at most $s_g\le d-1$ terms of $V$.
Consequently $V$ is regular over $pH$, and
Lemma~\ref{lem:c0-regular-sums} yields $\sum_0(V)=pH$.
The terms of $V$ are sums of pairwise disjoint subsequences of
$f^{ps_f}g_0^{ps_g}$, leaving $p$ occurrences of $f$ unused.
Since $\sum_0(f^p)+pH=H$, we obtain
\[
 H=\sum_0(V)+\sum_0(f^p)
 \subseteq\sum_0(f^{ps_f}g_0^{ps_g})
 \subseteq\sum_0(f^cg_0^b).
\]
Choose $T_1\mid f^cg_0^b$ with $\sigma(T_1)=-ah$.
Then $T=T_1h^a$ is nonempty and zero-sum, with $\mathsf v_h(T)=a=\mathsf v_h(S)$.
\end{proof}

\begin{proof}[Proof of Theorem~\ref{thm1.5}]
The lower bound for $n\ge2$ follows from Theorem~\ref{thm1.4}.
For $n=1$, both stated upper bounds follow from the known value of
$\mathsf N(C_p^2)$.
Fix $p$ and use strong induction on the integers $n$ whose prime divisors are at least $p$,
with the known value of $\mathsf N(C_p^2)$ as the base $n=1$.
Let $S$ be weak-regular over $G$, with $|S|\ge \operatorname{ub}(G)$.
If $\mathsf v_g(S)=\ord(g)$ for some $g\mid S$, then
$T=g^{\ord(g)}$ is zero-sum and $\mathsf v_g(T)=\mathsf v_g(S)$.
Thus every multiplicity is strictly smaller than the element order;
in particular, $0\nmid S$.
If $S$ is regular, Corollary~\ref{cor:regular-main} applies.

Otherwise choose $H$ by Lemma~\ref{lem:4.0}.
Then $H$ is nontrivial, $S_H$ is regular over $H$, and $|S_H|\ge|H|$.
First suppose $H=C_p\oplus C_{mp}$ is noncyclic.
If $m=1$, the known value of $\mathsf N(C_p^2)$ is at most $p^2\le|S_H|$.
If $m\ge2$ and $p\ge5$, then $|S_H|\ge mp^2\ge \operatorname{ub}(H)$,
so Corollary~\ref{cor:regular-main}, applied within $H$, suffices.
For $p=3$, use Lemma~\ref{lem:N-noncyclic-subgroup-three}.
In each case, $\mathsf v_h(S_H)=\mathsf v_h(S)$ for every $h\in H$,
so the subsequence obtained in $S_H$ also has the required property in $S$.

We may therefore suppose $H=C_m$ is cyclic.
If $|S_H|\ge\mathsf N(H)$, apply its definition within $H$.
Otherwise $m\le|S_H|\le\mathsf N(H)-1$.
Let $S_1=SS_H^{-1}$, and let $\Phi:G\to G_1=G/H$ be the quotient map.
If $G_1$ is cyclic, apply Lemma~\ref{lem:N-index-p-complete} when $|G_1|=p$
and Lemma~\ref{lem:N-cyclic-quotient} otherwise.
These lemmas apply because $|S|\ge \operatorname{ub}(G)\ge2np-1$.

Suppose finally that $G_1$ is noncyclic.
Since $G_1/pG_1\cong G/(H+pG)$ has order $p^2=|G/pG|$,
we obtain $H\le pG$.
Thus $1<m\mid n$ and $G_1\cong C_p\oplus C_{kp}$, where $k=n/m<n$.
If $\mathsf v_{\bar g}(\Phi(S_1))\ge\ord(\bar g)$ for some
$\bar g\in G_1\setminus\{0\}$ and
$|\Phi^{-1}(\bar g)\cap\supp(S_1)|\ge2$,
Lemma~\ref{lem:N-heavy-fiber} applies.
Otherwise, for each $g\mid S_1$ with
$\mathsf v_g(S_1)\ge r_g=\ord(\Phi(g))$, choose $\lfloor\mathsf v_g(S_1)/r_g\rfloor$ disjoint copies of
$g^{r_g}$.
Let $T_0$ be their product, $U$ the sequence of their sums,
and $R=S_1T_0^{-1}$ the remaining sequence.
Every term of $U$ is nonzero because $r_g\le\mathsf v_g(S)<\ord(g)$,
and $\Phi(R)$ is weak-regular.
If $|U|\ge m-1$, Lemma~\ref{lem:N-kernel-blocks} applies with $K=H$.
Otherwise $|T_0|\le(m-1)kp$, and
\[
 |R|\ge2np-1-(\mathsf N(H)-1)-(m-1)kp\ge\mathsf N(G_1).
\]
The last inequality is Lemma~\ref{lem2.7}, using the appropriate induction
bound for $k\ge2$, or the known elementary-group value for $k=1$.
By the definition of $\mathsf N(G_1)$, there is $R_1\mid R$ with
$\sigma(\Phi(R_1))=0$ and
$\mathsf v_{\bar g}(\Phi(R_1))=\mathsf v_{\bar g}(\Phi(R))$
for some $\bar g\mid\Phi(R)$.
Then $T_1=R_1T_0$ has sum in $H$ and contains all terms of $S_1$
whose image is $\bar g$. In particular,
$\mathsf v_h(T_1)=\mathsf v_h(S_1)$ for any
$h\in\Phi^{-1}(\bar g)\cap\supp(S_1)$.
Since $\sum_0(S_H)=H$, Lemma~\ref{lem2.9} completes the proof.
Replacing $n$ by $np$ proves the asserted equality for $C_p\oplus C_{np^2}$.
\end{proof}

\section{Remarks and open questions}
We conjecture that $\mathsf N(C_p\oplus C_{np})=2np-1$ whenever
$n\ge2$ and $p\ge3$ is the smallest prime dividing $np$.
The non-regular case in the proof of Theorem~\ref{thm1.5} uses only
$|S|\ge2np-1$. Hence, for the remaining cases $p\nmid n$, the bound
$c_0(C_p\oplus C_{np})\le\lceil(4np-2)/3\rceil$ would imply the
conjecture by Proposition~\ref{prop:regular-N-reduction}.

For the cyclic groups below, let $A_{n+1}=\lceil2\sqrt{n+1}\rceil$ and
\[
B_{n+1}=\min\{a+b: a,b\in\N,\ ab\ge n+1,\ \gcd(a,n)=1\},
\]
as in \cite{[JY2024]}.
For the following two families, $\mathsf N(C_n)$ attains the upper bound
$2n-A_{n+1}+1$ and is strictly larger than $2n-B_{n+1}+1$.

\begin{thm}\label{thm5.3}
	Let $n=2730\bigl(1+105(2t-1)\,(13t-\tfrac{13\pm 1}{2})\bigr)$ for a positive integer~$t$;
	equivalently, $n\in T_2\cup T_3$, where
	\begin{align*}
	T_2 &= \{\,2730(2730t^2-2835t+736)\,\}_{t\ge 1},\\
	T_3 &= \{\,2730(2730t^2-2625t+631)\,\}_{t\ge 1}.
	\end{align*}
	Then
	\[
	B_{n+1}=A_{n+1}+1=\lceil 2\sqrt{n+1}\rceil+1,
	\]
	and
	\[
	\mathsf{N}(C_n)=2n-B_{n+1}+2=2n-\lceil 2\sqrt{n+1}\rceil+1.
	\]
\end{thm}

\begin{proof}
By \cite[Lemmas~5.2 and~4.3]{[JY2024]},
$B_{n+1}=A_{n+1}+1$ and $\mathsf N(C_n)\le2n-A_{n+1}+1$.
Let $C_n=\langle e\rangle$.

Suppose first that $n\in T_2$, and let
$d=1365(2t-1)-52$ and $x=(12n-d)/13$.
Then $n=d(d-1)-26$, $13\mid d$, and $d$ is odd.
Since $d\ge1313$, we have
$(d-1)^2<n+1<(d-\tfrac12)^2$, and hence $A_{n+1}=2d-1$.
The identity $n=d(d-1)-26$ gives
\[
13x\equiv-d\pmod n,\qquad (d-1)x\equiv-2\pmod n.
\]
Every common divisor of $x$ and $n$ therefore divides both $d$ and $2$,
so $\gcd(x,n)=1$.
As $1<x<n$, the sequence $S=e^{n-d}(xe)^{n-d+1}$ is weak-regular
and has length $2n-A_{n+1}$.
If a zero-sum subsequence contains all $n-d+1$ copies of $xe$,
its multiplicity of $e$ is congruent to $(d-1)x\equiv-2\pmod n$.
It must therefore equal $n-2>n-d$, a contradiction.
If a zero-sum subsequence contains all $n-d$ copies of $e$,
then $13x\equiv-d\pmod n$ and $\gcd(x,n)=1$ force its
multiplicity of $xe$ to equal $n-13>n-d+1$, also a contradiction.

Suppose now that $n\in T_3$, and let
$d=1365(2t-1)+53$ and $x=(n-d)/2$.
Then $n=d(d-1)-26$, $d$ is even, and $d\equiv1\pmod{13}$.
In particular, $d\ge1418$ and $A_{n+1}=2d-1$ as above. We have
\[
2x\equiv-d\pmod n,\qquad (d-1)x\equiv-13\pmod n.
\]
Every common divisor of $x$ and $n$ divides both $d$ and $13$,
so $\gcd(x,n)=1$.
Again $1<x<n$, and $S=e^{n-d}(xe)^{n-d+1}$ is weak-regular
with $|S|=2n-A_{n+1}$.
A zero-sum subsequence containing all copies of $xe$ would have
$n-13>n-d$ copies of $e$, by the second congruence.
A zero-sum subsequence containing all copies of $e$ would have
$n-2>n-d+1$ copies of $xe$, by the first congruence and $\gcd(x,n)=1$.
Thus $S$ has no zero-sum subsequence $T$ with
$\mathsf v_g(T)=\mathsf v_g(S)$ for $g\in\{e,xe\}$.
In both cases, the definition of $\mathsf N(C_n)$ gives
$\mathsf N(C_n)\ge2n-A_{n+1}+1$, proving the equality.
\end{proof}

\begin{rem}
	The construction of Theorem~\ref{thm5.3} relies on the specific relation
	$n=d(d-1)-26$ together with the arithmetic of $13$ modulo $d$ (or $d-1$).
	For the other two families $T_4,T_5$ identified in \cite[Lemma~5.2]{[JY2024]}
	(which satisfy $n=d^2-64$ rather than $n=d(d-1)-26$), the analogous construction
	does not apply directly, and whether $\mathsf{N}(C_n)$ on $T_4\cup T_5$ attains
	the upper or the lower bound remains an open question.
\end{rem}

\bigskip
\noindent
{\bf Acknowledgement.}
The author thanks W. Gao for his guidance and helpful suggestions.
This work was supported in part by the National Natural Science Foundation
of China (Grant No. 12071344).
This paper is based on parts of Chapters~4 and~5 of the author's doctoral
dissertation, completed in 2025.

\bibliographystyle{amsplain}
\begin{thebibliography}{99}

\bibitem{[GG2006]} W. Gao and A. Geroldinger,
\emph{Zero-sum problems in finite abelian groups: a survey},
Expo. Math. \textbf{24} (2006), 337--369.

\bibitem{[GHQQZ2015]} W. Gao, D. Han, G. Qian, Y. Qu and H. Zhang,
\emph{On additive bases II}, Acta Arith. \textbf{168} (2015), 247--267.

\bibitem{[GHHLYZ2022-1]} W. Gao, S. Hong, W. Hui, X. Li, Q. Yin and P. Zhao,
\emph{Representation of zero-sum invariants by sets of zero-sum
sequences over a finite abelian group},
Period. Math. Hung. \textbf{85} (2022), 52--71.

\bibitem{[GHHLYZ2022-2]} W. Gao, S. Hong, W. Hui, X. Li, Q. Yin and P. Zhao,
\emph{Representation of zero-sum invariants by sets of zero-sum
sequences over a finite abelian group II},
J. Number Theory \textbf{241} (2022), 738--760.

\bibitem{[GLQW2022]} W. Gao, Y. Li, Y. Qu and Q. Wang,
\emph{Additive bases of finite abelian groups of rank 2},
Integers \textbf{22} (2022), Paper No. A37.

\bibitem{[GQZ2020]} W. Gao, Y. Qu and H. Zhang,
\emph{On additive bases III}, Acta Arith. \textbf{193} (2020), 293--308.

\bibitem{[GH2006]} A. Geroldinger and F. Halter-Koch,
\emph{Non-Unique Factorizations. Algebraic, Combinatorial and Analytic Theory},
Pure and Applied Mathematics, Vol.~278, Chapman \& Hall/CRC,
Boca Raton, FL, 2006.

\bibitem{[JY2024]} X. Jiang and W. Yang,
\emph{A zero-sum problem related to the max gap of the unit group of the residue class ring},
J. Number Theory \textbf{260} (2024), 198--211.

\bibitem{[N1996]} M. B. Nathanson,
\emph{Additive Number Theory: Inverse Problems and the Geometry of Sumsets},
Graduate Texts in Mathematics, Vol.~165, Springer, New York, 1996.

\bibitem{[QH2017]} Y. Qu and D. Han,
\emph{Additive bases of $C_p\oplus C_{p^n}$},
Int. J. Number Theory \textbf{13} (2017), 2453--2459.

\bibitem{[QL2022]} Y. Qu and Y. Li,
\emph{Additive bases of $C_3\oplus C_{3q}$},
Colloq. Math. \textbf{169} (2022), 189--195.

\bibitem{[QLWZ2025]} Y. Qu, Y. Li, Q. Wang and X. Zhao,
\emph{A note on additive bases of abelian groups of rank 2},
Graphs Combin. \textbf{41} (2025), Paper No. 24.

\end{thebibliography}

\end{document}
